\chapter{Extensiones supervivientes, monodromía y dinámica del cociente}
\label{chap:extensiones-dinamica-cociente}
\index{supervivencia!extensión global}
\index{monodromía!retorno nonádico}
\index{medida!catálogo de emisiones}

El lenguaje de los caminos necesita una precisión decisiva. Un camino de
APP, una emisión visible, una historia enriquecida de TPK y una sección
global no son cuatro nombres para el mismo objeto. Constituyen cuatro niveles
consecutivos de una misma construcción. El camino conserva el orden
posicional; la emisión conserva una lectura finita; la historia añade firma,
acarreo, orientación, frontera y memoria; la sección global exige que todos
sus prefijos sobrevivan simultáneamente a cualquier profundidad.

Esta jerarquía permite separar dos cuestiones que habían permanecido
entremezcladas. La primera es lógica: bajo qué hipótesis un prefijo finito
pertenece de verdad a una extensión infinita compatible. La segunda es
dinámica: qué significa transportar una sola historia y qué significa
transferir simultáneamente todas sus continuaciones compatibles. La
distinción es esencial en APP\@. Una trayectoria determinada lleva un estado a
otro; una transferencia de caminos lleva una distribución de extensiones a
otra distribución.

En esta sección se resuelve el problema finito de profundidad seis. La
cronología determinista publicada posee retorno visible y no puede seleccionar
una probabilidad única. En cambio, el transporte sesgado por la fase nonádica,
cuando cada canal activo se interpreta como la renovación normalizada de toda
su fibra de continuaciones compatibles, tiene un retorno de nueve pasos
uniforme sobre las (468) emisiones. Los pesos (1/3) dejan entonces de ser
una elección exterior: proceden de las tres apariciones de cada firma
(+1,-1,0) en el ciclo nonádico orientado. La normalización dentro de cada fibra sigue
siendo una premisa matemática precisa ---el conteo uniforme de todos los
caminos compatibles--- y no debe confundirse con el sucesor de una sola
trayectoria. El paso al límite inverso de todas las profundidades conserva un
enunciado propio, formulado al final de la sección.

\section{\(5\) niveles de la noción de camino}

Sea \(\Gamma_{\rm APP}\) la categoría de caminos cardinales del toro APP\@.
Para cada profundidad \(m\), sea \(\mathcal H_m\) el conjunto finito de
historias enriquecidas admisibles
\[
 h_m=(v_0,e_1,v_1,\ldots,e_m,v_m),
\]
donde cada flecha \(e_j\) transporta todos los campos declarados por TPK\@. Si
\(n\geq m\), la aplicación
\[
 p_{n,m}:\mathcal H_n\longrightarrow\mathcal H_m
\]
elimina la cola posterior a la profundidad \(m\). La sombra posicional
\[
 \operatorname{sh}_m:\mathcal H_m\longrightarrow\Gamma_{\rm APP}
\]
olvida la fibra de memoria y conserva el camino APP subyacente.

\begin{definition}[Jerarquía de extensiones]
\label{def:jerarquia-extensiones}
Distinguimos los objetos siguientes.
\begin{enumerate}
  \item Un \emph{camino APP} es un morfismo de
  \(\Gamma_{\rm APP}\).
  \item Una \emph{ruta observable} es una lectura de ese morfismo que
  conserva conjuntamente las hojas de suma y producto.
  \item Una \emph{extensión TPK} es una elevación de la ruta observable a
  una historia \(h_m\in\mathcal H_m\), con sus datos de residuo, cociente,
  firma, orientación, frontera, acarreo y registro.
  \item Una \emph{extensión superviviente} es una historia finita que admite
  prolongaciones compatibles a todo horizonte.
  \item Una \emph{sección global} es una familia compatible
  \(x=(h_m)_{m\geq0}\) de extensiones supervivientes.
\end{enumerate}
En particular, la ruta es la sombra de la extensión; no es un objeto exterior
que la extensión elija después.
\end{definition}

Para \(h\in\mathcal H_m\), su cilindro de prolongaciones es
\[
 Z_m(h)=
 \{x=(h_n)_n\in\varprojlim_n\mathcal H_n:x_m=h\}.
\]
En particular, toda familia de \(Z_m(h)\) satisface además
\(p_{n,\ell}(h_n)=h_\ell\) para \(n\geq\ell\); no se admiten colecciones de
truncamientos incompatibles.
Definamos también
\[
 \mathcal S_m^{(N)}=p_{N,m}(\mathcal H_N),
 \qquad
 \mathcal S_m^\infty=\bigcap_{N\geq m}\mathcal S_m^{(N)}.
\]

\begin{theorem}[Criterio condicional de extensión superviviente]
\label{thm:criterio-extension-superviviente}
Fijemos \(h\in\mathcal H_m\) y, para \(N\geq m\), escribamos
\[
 \mathcal P_N(h)=\{g\in\mathcal H_N:p_{N,m}(g)=h\}.
\]
Supóngase que los truncamientos son coherentes y que cada
\(\mathcal P_N(h)\) es finito. Son equivalentes:
\begin{enumerate}
  \item \(\mathcal P_N(h)\ne\varnothing\) para todo \(N\geq m\), es decir,
  \(h\in\mathcal S_m^\infty\);
  \item el árbol de prolongaciones finitas de \(h\) es finitamente ramificado
  y tiene un nivel no vacío a toda profundidad;
  \item \(Z_m(h)\ne\varnothing\).
\end{enumerate}
En consecuencia, siempre que los truncamientos supervivientes finitos sean no
vacíos, el conjunto
\[
 \Omega_{\rm H}^{\rm surv}
 =\varprojlim_m\mathcal S_m^\infty
\]
es exactamente el espacio de secciones globales supervivientes. El teorema no
afirma que esos truncamientos sean no vacíos para toda entrada APP: esa
propiedad debe demostrarse mediante el protocolo de supervivencia.
\end{theorem}

\begin{proof}
La implicación de la tercera condición a las dos primeras se obtiene por
truncamiento. Supongamos ahora la primera condición y ordenemos los elementos de
\(\bigsqcup_{N\geq m}\mathcal P_N(h)\) por la relación de prefijo. Cada nivel
es finito y no vacío, y la coherencia de los truncamientos hace de esta unión
un árbol finitamente ramificado con nodos a cualquier profundidad; esto da la
segunda condición. Recíprocamente, desde la segunda condición el lema de König
proporciona una rama infinita. Sus truncamientos forman una familia compatible
perteneciente a \(Z_m(h)\), que es la tercera condición. La misma construcción
identifica el límite inverso indicado allí donde sus truncamientos finitos son
no vacíos.
\end{proof}

El teorema resuelve la posible discordancia entre «prolongable hasta cada
horizonte por separado» y «posee una prolongación global»: la finitud de las
ramas es la hipótesis que permite escoger una única cadena compatible. No
afirma que todos los caminos APP pertenezcan a
\(\operatorname{sh}(\Omega_{\rm H}^{\rm surv})\). Esa cobertura exige probar
que los cierres enriquecidos no vacían las fibras correspondientes.

\section{Retorno de fase y relación de monodromía}

El retorno del cociente fase--longitud se formaliza por separado en
\cref{def:odometro-9-12,prop:proyeccion-fase-longitud}:
\[
(b,r)\xrightarrow{\;9\;}(b+1,r),
\qquad
(b,r)\xrightarrow{\;108\;}(b,r).
\]
Estas igualdades pertenecen al odómetro \(9\times12\). No constituyen la
prueba de \(F_{\rm obs}^{108}=I\), que usa el estado observable y aparece más
adelante.

Sea \(\phi\) la coordenada de fase con valores en \(\mathbb Z/9\mathbb Z\).
Un retorno visible satisface
\[
 \phi(h_{k+9})=\phi(h_k).
\]
El registro completo puede, sin embargo, cumplir
\[
 B(h_{k+9})\ne B(h_k).
\]
Esta desigualdad no es un fallo de periodicidad: es la memoria de la vuelta.
Una acción ordinaria de \(C_9\) sobre el estado completo impondría que la
novena potencia fuese la identidad y borraría precisamente esa información.

Para tipar el retorno, sea \(H\) una hexada y
\(P_2(H)\) su conjunto de quince pares. Sean \(\mathcal B_k\) los estados de
registro que sobreviven en la fase \(k\), sea
\(\mathcal T_k\subseteq\mathcal B_k\times\mathcal B_{k+1}\) la relación de
transporte publicada y supongamos dado un codificador de fibra
\[
 c_k:\mathcal B_k\longrightarrow P_2(H).
\]
Sólo después de declarar esos codificadores se define la relación observable
\[
 R_k=(c_k\times c_{k+1})(\mathcal T_k)
 \subseteq P_2(H)\times P_2(H).
\]
La monodromía relacional de una vuelta es
\[
 \mathcal M_k
 =R_{k+8}\circ R_{k+7}\circ\cdots\circ R_k.
\]

\begin{proposition}[Clasificación del retorno nonádico dado el codificador]
\label{prop:clasificacion-retorno-nonadico}
Se dan tres niveles lógicamente distintos.
\begin{enumerate}
  \item Sin una prueba de funcionalidad, el retorno es la relación
  \(\mathcal M_k\).
  \item Si cada \(R_k\) es la gráfica de una aplicación
  \(\mu_k:P_2(H)\to P_2(H)\), el transporte es el producto sesgado
  \[
   T(k,p)=(k+1,\mu_k(p)),
  \]
  y
  \[
   T^9(k,p)=(k,M_k(p)),
   \qquad
   M_k=\mu_{k+8}\circ\cdots\circ\mu_k.
  \]
  \item Si las \(\mu_k\) son biyectivas, \(M_k\) es una holonomía
  automorfa. El transporte desciende a una acción ordinaria de
  \(C_9\) sobre el estado completo si y sólo si \(M_k=I\) en todas las
  fibras.
\end{enumerate}
\end{proposition}

\begin{proof}
La primera afirmación es la definición de imagen y composición relacional una
vez fijados \((c_k)_k\). Bajo
funcionalidad, nueve aplicaciones consecutivas mantienen fija la primera
coordenada y componen sus acciones sobre la segunda. La composición de
biyecciones es una biyección. Finalmente, \(T^9=I\) equivale exactamente a
\(M_k=I\) para cada fase y cada punto de la fibra.
\end{proof}

\begin{corollary}[No retorno puntual bajo transporte acumulado]
\label{pf84:E02-29}
Sean \(q:X\to V\), \(T:X\to X\) y \(x\in X\). Si
\[
 q(T^9x)=q(x),\qquad
 T^9x=\mathcal M_x\,x,\qquad
 \mathcal M_x\,x\ne x,
\]
entonces la proyección retorna y el estado no:
\[
 q(T^9x)=q(x),\qquad T^9x\ne x.
\]
La hipótesis puntual \(\mathcal M_x\,x\ne x\) es indispensable; la desigualdad
global \(\mathcal M_x\ne I\) no excluye que \(x\) sea un punto fijo.
\end{corollary}

\begin{proof}
La primera igualdad es una hipótesis. La segunda y la separación puntual dan
\(T^9x=\mathcal M_x\,x\ne x\). El enunciado quedaría refutado por un caso que
satisfaga las tres premisas y tuviera \(T^9x=x\); no autoriza a sustituir la
tercera premisa por \(\mathcal M_x\ne I\).
\end{proof}

\begin{theorem}[Suficiencia funcional de la tupla enriquecida]
\label{pf84:E02-30}
Sea
\[
 E:\mathscr X_{\rm TPK}^{\rm enr}\longrightarrow\mathscr S_E
\]
la tupla que conserva bloque visible, bloques nuevos y acarreo, estado de
elevación, cilindros, firma de frontera, orientación, observable, fase y
registro dodecafásico. Sean
\[
 T:\mathscr X_{\rm TPK}^{\rm enr}\to
 \mathscr X_{\rm TPK}^{\rm enr},
 \qquad
 O:\mathscr X_{\rm TPK}^{\rm enr}\to Y.
\]
Cuando existen mapas
\[
 \widehat T:\mathscr S_E\to\mathscr X_{\rm TPK}^{\rm enr},
 \qquad
 \widehat O:\mathscr S_E\to Y
\]
tales que
\[
 T=\widehat T\circ E,\qquad O=\widehat O\circ E,
\]
la igualdad \(E(x)=E(y)\) implica \(T(x)=T(y)\), \(O(x)=O(y)\) y, por
determinismo, las mismas continuaciones y salidas posteriores.
\end{theorem}

\begin{proof}
Por factorización,
\[
 T(x)=\widehat T(E(x))=\widehat T(E(y))=T(y),
 \qquad
 O(x)=\widehat O(E(x))=\widehat O(E(y))=O(y).
\]
La primera igualdad permite iterar el argumento. Ésta es una condición
suficiente, no una equivalencia no probada. Dos estados con la misma tupla y
transiciones o salidas diferentes, el acceso a una coordenada ausente o la
consulta de un valor objetivo son falsadores. La microfibra visible de
cardinal cinco del \cref{pf84:E02-22} sólo controla la no fidelidad de la
palabra \(w_6\); no demuestra esta factorización positiva.
\end{proof}

La ventana finita publicada contiene nueve retornos de fase con cambio del
estado enriquecido y una selección local de tres continuaciones a una al
aumentar dos horizontes. Eso prueba retorno con memoria y refuta la
identificación del registro con un ciclo sin memoria. La extracción global de
los nueve codificadores \(c_k\) y, a través de ellos, de las relaciones
\(R_k\), desde todas las historias supervivientes es una operación distinta.
La proposición anterior clasifica el retorno una vez dado el mapa
registro--pares; no deriva ese mapa desde el estado enriquecido de TPK.

El episodio concreto de la novena fase, con sus tres continuaciones, el
bloque \(B_{11}\), la frontera \((502,662,638)\) y las desigualdades enteras
de cilindros, se demuestra una sola vez en el
\cref{thm:seleccion-novena-fase}. En esta residencia sólo se usa su
consecuencia tipada:
\[
\#\operatorname{Surv}_{1}(B_9)=3,\qquad
\#\operatorname{Surv}_{2}(B_9)=1.
\]

Las otras dos palabras pueden preceder la misma cola ternaria sólo al cambiar
la extensión profunda de Hensel. Comparten una parte del residuo visible,
pero no la genealogía del estado.

\section{El catálogo de 468 emisiones como producto de canales}

El catálogo finito parte del conjunto
\(\Omega_{\rm seed}\) de \(104\,976\) presentaciones microscópicas. Sea
\(\mathcal U_6\) el conjunto de emisiones séxtuples distintas y sea
\[
 F:\Omega_{\rm seed}\twoheadrightarrow\mathcal U_6.
\]
El censo exacto es
\[
 |\mathcal U_6|=468,
 \qquad
 \#\{U:|F^{-1}(U)|=144,432,1944\}=432,18,18.
\]
Estos conteos pertenecen a \(F\). Compararlos con historias TPK de
profundidad seis requeriría un mapa explícito desde el dominio de historias
hacia \(\Omega_{\rm seed}\); ningún resultado de esta sección presupone ese
mapa ni identifica presentaciones del catálogo con estados enriquecidos.
El catálogo finito proporciona, para cada emisión, una firma aditiva
\(s\) y una firma multiplicativa \(e\). La aplicación conjunta es una
biyección
\[
 \mathcal U_6\xrightarrow{\sim}
 \mathcal S_+\times\mathcal S_\times,
 \qquad
 |\mathcal S_+|=18,
 \quad
 |\mathcal S_\times|=26.
\]
Así,
\[
 \boxed{468=18\cdot26}
\]
es una factorización del objeto y no sólo de su cardinal. La primera
coordenada conserva cómo suma la emisión; la segunda, cómo acumula el
producto. Esta separación es la forma finita explícita de la
incompatibilidad suma--producto: ambas lecturas coexisten en el estado, pero
cada proyección olvida la firma complementaria.

La primera coordenada admite una parametrización adicional que será útil al
compararla con la frontera nonádica. Si una semilla aditiva tiene coordenadas
\((i,j)\in(\mathbb Z/9\mathbb Z)^2\) y orientación
\(\varepsilon=+1\) para las direcciones este--sur y \(-1\) para
norte--oeste, su firma completa queda determinada biyectivamente por
\[
 \bigl(i+j\bmod9,\varepsilon\bigr)
 \in(\mathbb Z/9\mathbb Z)\times\{\pm1\}.
\]
Por tanto,
\[
 \boxed{\mathcal S_+\simeq
 (\mathbb Z/9\mathbb Z)\times\{\pm1\}.}
\]
La coordenada producto no posee una uniformidad análoga: sus \(26\) firmas
tienen fibras de semillas de tamaños \(8\), \(24\) y \(108\), con histograma
\(24\times8+1\times24+1\times108=324\). Esta asimetría es estructural. El
factor \(18\) ya contiene fase y orientación; cualquier descenso posterior
del factor producto a una clase de pares debe conservar memoria adicional o
aceptar fibras no uniformes.

Definamos sobre \(\ell^2(\mathcal S_+\times\mathcal S_\times)\) las
esperanzas condicionales
\[
\begin{aligned}
 (E_+f)(s,e)&=\frac1{18}\sum_{s'\in\mathcal S_+}f(s',e),\\
 (E_\times f)(s,e)&=\frac1{26}\sum_{e'\in\mathcal S_\times}f(s,e').
\end{aligned}
\]
El primer operador renueva la firma aditiva y conserva la multiplicativa;
el segundo realiza la operación dual. El tercer canal retiene ambas.

\section{La trayectoria única y la transferencia de caminos}

A diferencia del retorno del odómetro definido en
\cref{def:odometro-9-12,prop:proyeccion-fase-longitud}, la proposición
siguiente demuestra el retorno del operador observable \(F_{\rm obs}\).

La ley observable publicada actualiza un solo cursor en cada paso. En las
fases \(+1\) avanza el cursor aditivo; en las fases \(-1\), el
multiplicativo; en las fases \(0\), ambos permanecen quietos. Después de los
pasos \(27,54,81,108\) se invierten las dos orientaciones. Denotemos esta
evolución determinista por \(F_{\rm obs}\).

\begin{proposition}[Retorno de la cronología observable]
\label{prop:retorno-fobs-identidad}
Para toda elección de las dos posiciones y orientaciones iniciales, las
posiciones y las orientaciones se restauran después de \(54\) pasos. Al
restaurar además la fase se obtiene
\[
 F_{\rm obs}^{108}=I.
\]
Por tanto, el retorno de Poincaré de \(108\) pasos induce la identidad sobre
\(\mathcal U_6\) y no selecciona una medida estacionaria única.
\end{proposition}

\begin{proof}
En cada bloque de \(27\) pasos, cada cursor realiza exactamente nueve
avances en una orientación fija. Como el espacio de posiciones es
\(X_9=(\mathbb Z/9\mathbb Z)^2\), su desplazamiento es \(9d=0\). Al final
del bloque se invierte \(d\). Dos bloques restauran posición y orientación;
cuatro restauran también el índice en \(\Theta_{108}\). El retorno inducido
es, por consiguiente, la identidad. Toda probabilidad es estacionaria para
la identidad, de modo que ninguna queda seleccionada por esa cronología.
\end{proof}

El problema aparece incluso antes del retorno si se intenta olvidar la
memoria de ruta. La firma producto
\[
 e_0=(5,5,5,5,5,5)
\]
tiene veinticuatro representantes. Tras un avance en su propia orientación,
esos representantes se separan, ocho a ocho, entre
\[
 (5,5,5,7,1,7),\qquad
 (5,5,5,7,7,1),\qquad
 (5,5,5,1,7,7).
\]
Así, el avance producto no desciende a una aplicación de las veintiséis
firmas: el representante que se había olvidado vuelve a ser relevante. Este
testigo explica por qué una ruta individual no puede convertirse, por simple
proyección, en el refresco de una fibra completa.

La lectura APP es otra. En lugar de escoger un representante, conserva todas
las continuaciones compatibles. Para formalizarla en el catálogo finito,
consideremos las subálgebras de funciones que sólo dependen de una de las dos
firmas. Respecto del conteo uniforme, \(E_+\) y \(E_\times\) son las
esperanzas condicionales que olvidan exactamente la coordenada activa. Son los
únicos proyectores de Markov que cumplen simultáneamente:
\begin{enumerate}
  \item preservación de la coordenada inactiva;
  \item olvido completo de la coordenada activa;
  \item preservación del estado de conteo uniforme.
\end{enumerate}
En efecto, una fila que ha olvidado la coordenada de partida debe ser
constante sobre los \(18\), respectivamente \(26\), destinos compatibles; la
normalización fuerza los valores \(1/18\) y \(1/26\). Ésta es la canonicidad
relativa de la renovación: no se postula una preferencia entre ramas que APP
declara igualmente posibles.

Sea \(C_9=\mathbb Z/9\mathbb Z\), con la palabra de fase
\[
 \tau=(+1,-1,0,+1,-1,0,+1,-1,0).
\]
Escribamos
\[
 P_{+1}=E_+,
 \qquad P_{-1}=E_\times,
 \qquad P_0=I.
\]

\begin{definition}[Transferencia TPK sesgada por la fase]
\label{def:transferencia-tpk-fase}
Sobre \(\mathcal U_6\times C_9\) definimos
\[
 \mathcal K_{\rm ph}\bigl((u,r),(v,r+1)\bigr)
 =P_{\tau(r)}(u,v),
\]
y hacemos cero las demás entradas. Esta definición es relativa al lector de
todos los caminos compatibles: una fase activa renueva uniformemente la fibra
que esa fase lee y la fase neutra retiene la emisión.
\end{definition}

\begin{theorem}[Retorno nonádico uniforme]
\label{thm:retorno-nonadico-uniforme}
El núcleo \(\mathcal K_{\rm ph}\) es doblemente estocástico e irreducible,
tiene período exactamente nueve y posee un único estado estacionario,
\[
 \pi_{\rm ph}(u,r)=\frac1{9\cdot468}.
\]
En cada fibra de fase, su retorno de nueve pasos es
\[
 \mathcal K_{\rm ph}^{\,9}\big|_r=E_+E_\times,
 \qquad
 (E_+E_\times)(u,v)=\frac1{468}.
\]
Si la fase inicial se distribuye uniformemente, la sombra de un paso sobre
las emisiones es
\[
 \overline K_{\rm ph}=\frac13(I+E_++E_\times).
\]
\end{theorem}

\begin{proof}
Cada \(P_{\tau(r)}\) es doblemente estocástico y el desplazamiento
\(r\mapsto r+1\) es una permutación; su producto sesgado conserva ambas sumas.
Los tres operadores son idempotentes y conmutan. Como cada uno aparece tres
veces en la palabra nonádica,
\[
 \prod_{r\in C_9}P_{\tau(r)}
 =E_+^3E_\times^3I^3=E_+E_\times.
\]
La primera esperanza olvida la firma aditiva y la segunda la multiplicativa,
de modo que su composición asigna a cada uno de los \(18\cdot26=468\)
destinos la probabilidad \(1/468\). Por ello, en nueve pasos se comunican dos
estados cualesquiera de una misma fase; los pasos restantes comunican las
fases, y el núcleo es irreducible. Todo retorno debe tener longitud múltiplo
de nueve, mientras que la diagonal de \(\mathcal K_{\rm ph}^9\) es positiva;
el período es nueve. La doble estocasticidad y la irreducibilidad dan el
estado uniforme único. Finalmente, promediar las nueve fases cuenta tres
copias de cada canal y produce la última fórmula.
\end{proof}

Olvidar la fase no es un cociente de Markov fuerte: desde una misma emisión,
dos fases distintas aplican \(I\), \(E_+\) o \(E_\times\), que tienen filas
distintas. El operador \(\overline K_{\rm ph}\) es una sombra promediada en la
fase estacionaria, no el sucesor cronológico de una emisión sin fase. Esta
distinción conserva a la vez la dinámica y el significado de la medida.

\begin{definition}[Núcleo promediado inducido por la transferencia]
\label{def:nucleo-tritico-refresco}
Denominamos núcleo promediado del catálogo a la sombra de fase
\[
 \boxed{
 K_{\rm cat}=\frac13(I+E_++E_\times).
 }
\]
Los pesos son inducidos por el calendario TPK, que contiene tres ocurrencias
de cada signo. Los operadores de renovación son inducidos relativamente a la
normalización uniforme de todas las continuaciones compatibles de la fibra
activa. Ninguna de las dos afirmaciones identifica \(K_{\rm cat}\) con
\(F_{\rm obs}\).
\end{definition}

\begin{proposition}[Familia convexa de núcleos estocásticos con medida uniforme común]
\label{prop:familia-refrescos-abc}
Para \(a,b,c\geq0\), \(a+b+c=1\), la familia
\[
 K_{a,b,c}=aI+bE_++cE_\times
\]
está formada por núcleos simétricos y doblemente estocásticos. Si \(b,c>0\),
son irreducibles y, por tener diagonal positiva, aperiódicos. Su espectro es
\[
 \operatorname{Spec}(K_{a,b,c})=
 \{1^{[1]},(a+c)^{[17]},(a+b)^{[25]},a^{[425]}\}.
\]
En particular, existe un continuo de dinámicas con la misma factorización y la
misma medida uniforme. La elección
\(K_{\rm cat}=K_{1/3,1/3,1/3}\) no queda caracterizada por el censo solo. La
seleccionan conjuntamente la palabra de fase TPK y la transferencia
normalizada de la \cref{def:transferencia-tpk-fase}.
\end{proposition}

\begin{proof}
Los tres sumandos son núcleos simétricos y estocásticos; por convexidad ocurre
lo mismo con \(K_{a,b,c}\), y la simetría convierte la estocasticidad por filas
en estocasticidad por columnas. Si \(b,c>0\), dos pasos permiten renovar
sucesivamente ambas coordenadas y comunicar cualquier par de estados. Además,
\[
 K_{a,b,c}((s,e),(s,e))=a+\frac b{18}+\frac c{26}>0,
\]
de modo que el núcleo es aperiódico. En la descomposición ortogonal
\[
 \mathbf1\oplus\ell_0^2(\mathcal S_+)
 \oplus\ell_0^2(\mathcal S_\times)
 \oplus\bigl(\ell_0^2(\mathcal S_+)\otimes
              \ell_0^2(\mathcal S_\times)\bigr),
\]
los pares \((E_+,E_\times)\) actúan, respectivamente, como
\((1,1),(0,1),(1,0),(0,0)\). Se obtienen así los cuatro autovalores y las
multiplicidades \(1,17,25,425\). Para \(b,c>0\), la irreducibilidad hace única
la medida estacionaria uniforme, lo que exhibe la no unicidad incluso bajo
esa exigencia.
\end{proof}

\begin{theorem}[Espectro exacto del refresco equiponderado]
\label{thm:dinamica-cociente-468}
El núcleo \(K_{\rm cat}\) es simétrico, doblemente estocástico,
irreducible y aperiódico. Su espectro, con multiplicidades, es
\[
 \boxed{
 \operatorname{Spec}(K_{\rm cat})
 =\left\{1^{[1]},
 \left(\frac23\right)^{[42]},
 \left(\frac13\right)^{[425]}\right\}.
 }
\]
En consecuencia, su brecha espectral es (1/3) y su único estado
estacionario es
\[
 \boxed{
 \tau_{468}(f)=\frac1{468}\sum_{U\in\mathcal U_6}f(U).
 }
\]
\end{theorem}

\begin{proof}
Los operadores \(E_+\) y \(E_\times\) son proyectores ortogonales de
promedio y conmutan. Las filas y columnas de cada uno suman uno; lo mismo
ocurre con su promedio. Desde \((s,e)\) se alcanza \((s',e')\) en dos
renovaciones, y la contribución de \(I\) proporciona una autoarista. Esto
prueba irreducibilidad y aperiodicidad.

El espacio se descompone ortogonalmente en cuatro sectores:
\[
 \mathbf1,
 \qquad
 \ell_0^2(\mathcal S_+),
 \qquad
 \ell_0^2(\mathcal S_\times),
 \qquad
 \ell_0^2(\mathcal S_+)\otimes
 \ell_0^2(\mathcal S_\times).
\]
Sus dimensiones son \(1,17,25,17\cdot25\). En el sector constante actúan
los tres sumandos como la identidad; en cada sector de una sola coordenada
actúan como (I+I+0); en el sector de interacción, como (I+0+0).
Dividir por tres da los autovalores y multiplicidades anunciados. Un núcleo
finito irreducible y aperiódico posee un único estado estacionario; como es
doblemente estocástico, éste es uniforme.
\end{proof}

El teorema es exacto relativamente a la factorización
\(\mathcal U_6\simeq\mathcal S_+\times\mathcal S_\times\), al calendario
TPK y al conteo uniforme de las continuaciones compatibles de cada fibra. La
familia \(K_{a,b,c}\) demuestra que la dinámica no se deduce del censo por sí
solo. La \cref{prop:retorno-fobs-identidad} distingue la transferencia de
todos los caminos de la cronología de uno de ellos.

\section{Elevación compensada construida desde el cociente}

Escribamos \(m(U)=|F^{-1}(U)|\). A partir del núcleo ya elegido en el cociente
definimos el núcleo de presentaciones por
\[
 \boxed{
 K_{\rm seed}(x,y)
 =\frac{K_{\rm cat}(F(x),F(y))}{m(F(y))}
 }
\]
Este núcleo es estocástico. Al sumar sobre una fibra de destino se obtiene
\[
 \sum_{F(y)=V}K_{\rm seed}(x,y)
 =K_{\rm cat}(F(x),V),
\]
independientemente del representante \(x\). Esta identidad muestra que el
elevación construida proyecta exactamente sobre \(K_{\rm cat}\). No demuestra que
una dinámica microscópica preexistente del estado enriquecido de TPK descienda
al catálogo.

\begin{theorem}[Elevación reversible construida desde el cociente]
\label{thm:elevacion-reversible-cociente}
La probabilidad
\[
 \widetilde\tau(x)
 =\frac1{468\,m(F(x))}
\]
es estacionaria y reversible para \(K_{\rm seed}\), y
\(F_\#\widetilde\tau=\tau_{468}\). Si
\[
 J_F\delta_U
 =m(U)^{-1/2}\sum_{F(x)=U}\delta_x,
 \qquad
 D=\operatorname{diag}(m(U)),
\]
entonces
\[
 K_{\rm seed}J_F
 =J_F\bigl(D^{1/2}K_{\rm cat}D^{-1/2}\bigr).
\]
En particular, el entrelazamiento hilbertiano sin conjugación por \(D\)
ocurre si y sólo si \(K_{\rm cat}\) conmuta con las multiplicidades de las
fibras.
\end{theorem}

\begin{proof}
La normalización se sigue de que cada una de las 468 fibras recibe masa
total (1/468). Usando la simetría de \(K_{\rm cat}\), para
\(F(x)=U\) y \(F(y)=V\) se tiene
\[
 \widetilde\tau(x)K_{\rm seed}(x,y)
 =\frac{K_{\rm cat}(U,V)}{468m(U)m(V)}
 =\widetilde\tau(y)K_{\rm seed}(y,x).
\]
Esto demuestra balance detallado, estacionariedad y la identidad de la
medida imagen.
Aplicar ambos miembros de la última identidad a un vector base
\(\delta_V\) da, sobre la fibra \(U\), el coeficiente
\(K_{\rm cat}(U,V)/\sqrt{m(V)}\); ése es exactamente el coeficiente del
operador conjugado mostrado.
\end{proof}

\begin{theorem}[Elevación de fase fiel a la quietud]
\label{thm:elevacion-fase-semillas}
Sea \(X=\Omega_{\rm seed}\), \(m(u)=|F^{-1}(u)|\), y definamos
\[
\begin{aligned}
 L_0(x,y)&=\mathbf1_{\{x=y\}},\\
 L_+(x,y)&=\frac{E_+(F(x),F(y))}{m(F(y))},\\
 L_\times(x,y)&=\frac{E_\times(F(x),F(y))}{m(F(y))}.
\end{aligned}
\]
Sobre \(X\times C_9\), sea
\[
 \widehat{\mathcal K}_{\rm ph}\bigl((x,r),(y,r+1)\bigr)
 =L_{\tau(r)}(x,y).
\]
Entonces \((F,\mathrm{id}_{C_9})\) satisface la propiedad de
\emph{agregabilidad fuerte} (\emph{strong lumpability}) de
\(\widehat{\mathcal K}_{\rm ph}\) hacia \(\mathcal K_{\rm ph}\). El núcleo
microscópico es irreducible, tiene período nueve y su estado estacionario
único es
\[
 \widetilde\pi_{\rm ph}(x,r)
 =\frac1{9\cdot468\,m(F(x))}.
\]
Su retorno nonádico satisface
\[
 \widehat{\mathcal K}_{\rm ph}^{\,9}
 \bigl((x,r),(y,r)\bigr)
 =\frac1{468\,m(F(y))}.
\]
\end{theorem}

\begin{proof}
La suma de \(L_+\), respectivamente \(L_\times\), sobre una fibra de destino
es \(E_+\), respectivamente \(E_\times\), e \(L_0\) proyecta a la identidad.
Esto prueba la agregabilidad en cada fase. Los factores \(m(F(y))\) cancelan
al componer: primero se selecciona uniformemente la firma activa y después un
representante de la emisión obtenida. Así, \(L_+\) y \(L_\times\) son
proyectores conmutantes y
\[
 L_+L_\times(x,y)=\frac1{468\,m(F(y))}.
\]
La positividad de esta entrada comunica todos los representantes en nueve
pasos. La fase fuerza que las longitudes de retorno sean múltiplos de nueve,
y la entrada diagonal del retorno es positiva. La fórmula estacionaria se
comprueba sumando primero los \(m(u)\) representantes de cada emisión; cada
una de las \(9\cdot468\) clases fase--emisión recibe la misma masa.
\end{proof}

Esta elevación mejora la lectura del núcleo promediado
\(K_{\rm seed}\): en la fase neutra conserva el representante microscópico,
en vez de remezclar silenciosamente su fibra. Al promediar y olvidar la fase,
ambas construcciones tienen la misma sombra \(K_{\rm cat}\), pero no la misma
dinámica oculta. La diferencia es necesaria para respetar que el cero de TRIT
es memoria de quietud.

\section{Congruencias mínimas estables bajo avance, media vuelta y 4.º de giro: \(28\) y \(162\) clases de memoria}

El testigo \(e_0=(5,5,5,5,5,5)\) muestra que la firma producto es demasiado
gruesa para transportar una ruta. Es posible calcular exactamente la memoria
mínima que falta. Sea
\[
 \mathcal X_\times=X_9\times\mathscr D_{\rm card}
\]
el espacio de \(324\) semillas producto. Denotemos por \(P\) el avance de una
casilla en la orientación almacenada, por \(H\) la media vuelta de esa
orientación y por \(Q\) el cuarto de giro
\[
 N\longmapsto E\longmapsto S\longmapsto O\longmapsto N.
\]

\begin{theorem}[Congruencias de avance y giro]
\label{thm:congruencias-avance-giro}
La congruencia mínima que refina las veintiséis firmas producto y es estable
por \(P\) y \(H\) posee \(28\) clases: veintisiete tienen tamaño ocho y una
tiene tamaño \(108\). La única escisión nueva es
\[
 (5,5,5,5,5,5)\longrightarrow
 \mathcal L_0\sqcup\mathcal L_1\sqcup\mathcal L_2,
 \qquad |\mathcal L_j|=8.
\]
Sus órbitas bajo \(P,H\) tienen tamaños
\[
 9,\quad9,\quad9,\quad1.
\]
Al combinar este refinamiento con las dieciocho firmas aditivas se obtiene
\[
 18\cdot28=504=468+36.
\]

La congruencia mínima estable por \(P\) y \(Q\) posee \(162\) clases de dos
semillas. Cada clase es una órbita de la involución
\[
 J(i,j,d)=\bigl(7-i,\,7-j,\,\overline d\bigr)\pmod9.
\]
Los operadores inducidos por \(P,Q\) actúan transitivamente sobre las
\(162\) clases.
\end{theorem}

\begin{proof}
Se parte de la partición por las veintiséis firmas y se refina
iterativamente una clase siempre que dos de sus elementos tengan imágenes en
clases distintas por alguno de los generadores. El proceso termina porque
\(\mathcal X_\times\) es finito. El censo exhaustivo produce, para
\((P,H)\), el histograma \(27\times8+1\times108=324\), las cuatro órbitas
indicadas y tres clases sobre la fibra excepcional de tamaño veinticuatro.
Para \((P,Q)\) produce \(162\times2=324\). La comprobación directa demuestra
que \(J\) preserva cada firma, conmuta con ambos generadores y empareja
exactamente los dos elementos de cada clase; el grafo de clases generado por
\(P,Q\) es conexo.
\end{proof}

El exceso \(36\) y el objeto \(R_{36}\) tienen tipos distintos. La frontera
de profundidad treinta y seis es
\[
 R_{36}=
 \begin{pmatrix}
 222220\\021222\\102011
 \end{pmatrix},
\]
es decir, tres bloques ordenados de seis trits, uno por canal. No es un
conjunto de treinta y seis estados. Por su parte, el término
\(504-468=36\) cuenta memoria adicional de hoja: la media vuelta que hace
funcional el transporte producto añade exactamente dos hojas nuevas a una
sola firma y, al repetirlas sobre las dieciocho firmas aditivas, produce
\(18(3-1)=36\). La coincidencia del entero (36) expresa dos graduaciones
del transporte, profundidad de frontera y multiplicidad de hoja, pero no
autoriza a identificar sus objetos.

El cuarto de giro \(Q\) tampoco estaba escondido en \(F_{\rm obs}\). La
cronología publicada contiene avances y medias vueltas programadas; incorporar
\(Q\) es una ampliación explícita del alfabeto de rutas. Una vez incorporado,
un paseo simétrico perezoso en \(P^{\pm1},Q^{\pm1}\) es irreducible y
aperiódico sobre las \(162\) clases, y su medida uniforme es única. Ésta es la
forma matemática exacta de «contar los giros»: el giro es una coordenada de
transporte y no una corrección decimal posterior.

\section{Cociclos de modo y orientación}
\label{sec:cociclos-modo-orientacion}

La misma regla permite definir los contadores operativos sin introducir la
Mecánica Dimensional. Añadamos al estado el último modo activo
\(m\in\{\Sigma,\Pi\}\). Una firma \(+1\) fija \(m=\Sigma\), una firma
\(-1\) fija \(m=\Pi\), y una firma \(0\) conserva el modo anterior. Para una
flecha \(a\), definamos el cociclo firmado
\[
 \omega_{\rm sw}(a)=
 \begin{cases}
  +1,&\Pi\longrightarrow\Sigma,\\
  -1,&\Sigma\longrightarrow\Pi,\\
  0,&\text{el modo se conserva}.
 \end{cases}
\]
Su variación positiva es \(s(a)=|\omega_{\rm sw}(a)|\). Definamos además
\(g(a)=1\) cuando cambia el par ordenado de orientaciones cardinales y
\(g(a)=0\) cuando se conserva. Al mantener en los extremos el modo y la
orientación terminal, las sumas
\[
 \Omega_{\rm sw}(\gamma)=\sum_{a\subset\gamma}\omega_{\rm sw}(a),
 \qquad
 S(\gamma)=\sum_{a\subset\gamma}s(a),
 \qquad
 G(\gamma)=\sum_{a\subset\gamma}g(a)
\]
son aditivas por concatenación.

En una vuelta nonádica cerrada, la sucesión de modos es
\[
 (\Sigma,\Pi,\Pi,\Sigma,\Pi,\Pi,\Sigma,\Pi,\Pi).
\]
Por tanto,
\[
 \Omega_{\rm sw}(C_9)=0,\qquad S(C_9)=6:
\]
hay tres cambios \(\Pi\to\Sigma\), tres cambios
\(\Sigma\to\Pi\) y tres retenciones. En el ciclo de \(108\) pasos,
\[
 \#(+1)=\#(-1)=\#(0)=36,\qquad
 \Omega_{\rm sw}=0,\qquad S=72,\qquad G=4.
\]
El último número cuenta las cuatro inversiones simultáneas programadas en
\(27,54,81,108\). Estos cociclos proceden directamente del calendario TPK\@.

\subsection{Cociente de modos y correspondencia areal--volumétrica}
\label{sec:cociente-modos-fav}

El signo del calendario, la etiqueta de modo y la hoja geométrica son tres
objetos tipados. En particular, no se identifica el signo de fase con el
observable angular sobre dígitos. Para la correspondencia geométrica fijamos el
lector HMT tipado
\[
 \mathfrak g_{\rm av}(\Sigma)=\mathscr H_{\rm ar},
 \qquad
 \mathfrak g_{\rm av}(\Pi)=\mathscr H_{\rm vol}.
\]
La primera asignación expresa la lectura aditiva o de borde; la segunda, la
lectura multiplicativa o de interior. El mapa \(\mathfrak g_{\rm av}\) es
parte del tipado de hojas. El resultado siguiente demuestra que, una vez
fijado ese tipado, la dinámica TPK fuerza la incidencia entre ellas.

\begin{theorem}[Descenso necesario de la incidencia de hojas]
\label{thm:descenso-necesario-fav}
Sea \(m_r\) el modo obtenido después de aplicar la fase \(r\) del bloque
primitivo \((+1,-1,0)\), con la regla
\[
 +1:\ m\mapsto\Sigma,\qquad
 -1:\ m\mapsto\Pi,\qquad
 0:\ m\mapsto m.
\]
La órbita cíclica de modos es \((\Sigma,\Pi,\Pi)\). Si
\[
 N_3(i,j)
 :=
 \#\{r\in\mathbb Z/3\mathbb Z:(m_r,m_{r+1})=(i,j)\},
\]
entonces, en el orden \((\Sigma,\Pi)\),
\[
 \boxed{
 N_3=
 \begin{pmatrix}0&1\\1&1\end{pmatrix}.}
\]
Por repetición del calendario,
\[
 N_9=3N_3,\qquad N_{27}=9N_3,\qquad N_{108}=36N_3.
\]
Después de aplicar \(\mathfrak g_{\rm av}\), la matriz primitiva de
incidencia es, en el orden
\((\mathscr H_{\rm ar},\mathscr H_{\rm vol})\),
\[
 \boxed{
 F_{\rm av}=
 \begin{pmatrix}0&1\\1&1\end{pmatrix}.}
\]
Así, las dos transiciones cruzadas, la única autorretención volumétrica, la
ausencia de autorretención areal y las multiplicidades primitivas unitarias
son consecuencias del cociente de modos de TPK; no son cuatro postulados
añadidos.
\end{theorem}

\begin{proof}
Los tres pares consecutivos, con cierre cíclico, son
\[
 \Sigma\longrightarrow\Pi,\qquad
 \Pi\longrightarrow\Pi,\qquad
 \Pi\longrightarrow\Sigma.
\]
Cada uno aparece una vez y el par
\(\Sigma\to\Sigma\) no aparece. Esto determina las cuatro entradas de
\(N_3\). Los calendarios de longitudes \(9,27,108\) contienen,
respectivamente, \(3,9,36\) copias del bloque primitivo, por lo que sus
matrices de incidencia son los múltiplos indicados. El máximo común divisor
de las entradas no nulas de cada matriz larga elimina únicamente esa
repetición global y devuelve \(N_3\). Finalmente,
\(\mathfrak g_{\rm av}\) transporta \(\Sigma,\Pi\) a las hojas areal y
volumétrica sin alterar la incidencia.
\end{proof}

La realización hiperbólica real que utiliza esta matriz y prueba su
invariancia lorentziana se formula en
\cref{sec:tres-generadores-concretos-trit}; allí \(F_{\rm av}\) se usa
como antecedente, sin repetir su derivación.

Este teorema debe distinguirse de una afirmación de Markov más fuerte.
Olvidar la fase no es un cociente de Markov fuerte de
\(\mathcal K_{\rm ph}\): el modo visible no determina cuál de los tres
operadores \(I,E_+,E_\times\) actúa en la iteración elemental siguiente. Lo que desciende
sin hipótesis adicional es el multigrafo de aristas de un bloque TPK\@.

Existe, sin embargo, una completación de memoria uno determinada de forma
única por ese descenso. Sea \(X_{\rm av}\) el menor subshift de tipo finito
sobre
\(\{\mathscr H_{\rm ar},\mathscr H_{\rm vol}\}\) que contiene todas las
palabras obtenidas al olvidar la fase y que sólo conserva pares consecutivos.
Todo subshift de un paso con esa propiedad debe admitir exactamente los tres
pares observados; el menor de ellos excluye el único par no observado. Por
tanto su matriz de adyacencia es \(F_{\rm av}\).

\begin{corollary}[Medida de Parry del envolvente de memoria uno]
\label{cor:parry-envolvente-fav}
El subshift \(X_{\rm av}\) es primitivo, tiene entropía topológica
\(\log\varphi\) y posee una única medida de máxima entropía. Sus pesos
estacionarios de vértice satisfacen
\[
 \boxed{
 \frac{\mu_{\rm ar}}{\mu_{\rm vol}}=\varphi^{-2}.}
\]
\end{corollary}

\begin{proof}
La raíz de Perron de \(F_{\rm av}\) es \(\varphi\), y un vector de Perron
positivo es \((1,\varphi)\). Como la matriz es simétrica, los pesos de vértice
de Parry son proporcionales al producto de los vectores izquierdo y derecho,
es decir, a \((1,\varphi^2)\).
\end{proof}

La medida de Parry pertenece al envolvente de caminos de memoria uno. No es
la imagen de \(\tau_{468}\), ni la frecuencia temporal de la órbita periódica
de modos. Esta última es
\[
 \nu_{\rm cyc}(\mathscr H_{\rm ar})=\frac13,\qquad
 \nu_{\rm cyc}(\mathscr H_{\rm vol})=\frac23,
\]
mientras \(\tau_{468}\) es uniforme sobre emisiones y la medida de Parry
pondera historias de renovación de longitud arbitraria. Las tres medidas
responden a preguntas distintas. La razón irracional
\(\varphi^{-2}\) queda así derivada canónicamente del envolvente de caminos,
sin atribuirla a un cociente de cardinales finitos.

\subsubsection{Del conteo cronológico al retorno marcado por \texorpdfstring{\(\pi\)}{π}}
\label{sec:retorno-marcado-pi-phi2}

Las dos fórmulas
\[
 \left(\frac13,\frac23\right)
 \qquad\text{y}\qquad
 \frac{\mu_{\rm ar}}{\mu_{\rm vol}}=\varphi^{-2}
\]
no son aproximaciones una de otra. La primera cuenta los tres instantes del
calendario primitivo \((\Sigma,\Pi,\Pi)\); la segunda pondera todas las
historias admisibles del envolvente de memoria uno. Por eso la frecuencia
cronológica es racional y la razón de Parry es irracional. En lo que sigue
escribimos
\[
 \frac{\pi_{\rm HMT}}{\varphi^2}
 =\pi_{\rm HMT}\varphi^{-2};
\]
no se trata de dividir por \(\varphi^{-2}\).

La aparición de \(\varphi^{-2}\) puede verse directamente en la dinámica
lineal. Si \(F_n\) denota el \(n\)-ésimo número de Fibonacci,
\[
 F_{\rm av}^{\,n}
 =
 \begin{pmatrix}
  F_{n-1}&F_n\\
  F_n&F_{n+1}
 \end{pmatrix}
 \qquad(n\geq1).
\]
Los autovalores de \(F_{\rm av}\) son
\(\varphi\) y \(-\varphi^{-1}\). La rama contraída invierte orientación en
una iteración. Al pasar a la memoria cuadrática, es decir, a
\(\operatorname{Sym}^2(\mathbb R^2)\), esa inversión se registra antes de
desaparecer el signo. En la base \((x^2,2xy,y^2)\), representando por
filas las imágenes de los vectores básicos, el operador inducido es
\[
 \operatorname{Sym}^2(F_{\rm av})=
 \begin{pmatrix}
  0&0&1\\
  0&1&2\\
  1&1&1
 \end{pmatrix},
 \qquad
 \chi(t)=(t+1)(t^2-3t+1).
\]
Sus tres autovalores son
\[
 \varphi^2,\qquad -1,\qquad \varphi^{-2}.
\]
Así, \(\varphi^{-2}\) es el modo estable positivo de la memoria de segundo
orden. Este paso explica por qué la razón áurea aparece aquí como ley de
retorno y no como media estadística añadida.

La clausura \(\pi\) pertenece a otro nivel del mismo estado HMT: es la salida
coinductiva del lector regional ya construido, no un autovalor nuevo de
\(F_{\rm av}\). Sean
\[
 P_{\rm ar}=
 \begin{pmatrix}1&0\\0&0\end{pmatrix},
 \qquad
 P_{\rm vol}=
 \begin{pmatrix}0&0\\0&1\end{pmatrix}.
\]
Entre los operadores positivos que preservan las dos hojas, hay uno y sólo
uno cuya normalización volumétrica es uno y cuya marca areal es la clausura
generada \(\pi_{\rm HMT}\):
\[
 D_\pi=\pi_{\rm HMT}P_{\rm ar}+P_{\rm vol}.
\]
La unicidad es relativa a estas tres cláusulas explícitas: positividad,
diagonalidad respecto de las hojas y las dos normalizaciones. Con
\(e_{\rm ar}=(1,0)^T\) y \(e_{\rm vol}=(0,1)^T\), el retorno marcado de
longitud \(n\) es
\[
 \Theta_n
 :=
 \frac{\langle e_{\rm ar},
 D_\pi F_{\rm av}^{\,n}e_{\rm ar}\rangle}
 {\langle e_{\rm vol},
 F_{\rm av}^{\,n}e_{\rm vol}\rangle}
 =
 \pi_{\rm HMT}\frac{F_{n-1}}{F_{n+1}}.
\]
Por tanto,
\[
 \boxed{\displaystyle
 \lim_{n\to\infty}\Theta_n
 =\frac{\pi_{\rm HMT}}{\varphi^2}.}
\]
La marca \(\pi_{\rm HMT}\) se aplica una sola vez, en el retorno de frontera;
no se reinyecta en cada transición. La combinatoria finita produce el modo
estable \(\varphi^{-2}\), el lector coinductivo produce la clausura
\(\pi_{\rm HMT}\), y \(D_\pi\) compone ambos datos sin invertir su orden
causal.

La longitud nativa \(108\) ofrece un testigo exacto de esa convergencia:
\[
 (F_{\rm av}^{108})_{\rm ar,ar}
 =F_{107}=10284720757613717413913,
\]
\[
 (F_{\rm av}^{108})_{\rm vol,vol}
 =F_{109}=26925748508234281076009.
\]
En consecuencia,
\[
 \left|
 \frac{F_{107}}{F_{109}}-\varphi^{-2}
 \right|
 =
 6.1684979916614407936\ldots\times10^{-46},
\]
y, después de la única marca de clausura,
\[
 \left|
 \pi_{\rm HMT}\frac{F_{107}}{F_{109}}
 -\frac{\pi_{\rm HMT}}{\varphi^2}
 \right|
 =
 1.9378907974286976068\ldots\times10^{-45}.
\]
Estos errores miden la aproximación finita del retorno; no son ajustes de
\(\pi\) ni de \(\varphi\).

\paragraph{Lector cilíndrico de la razón.}
Sean \(I_n^\pi=[\underline\pi_n,\overline\pi_n]\) e
\(I_n^\varphi=[\underline\varphi_n,\overline\varphi_n]\) los cilindros
positivos y encajados producidos por los dos lectores HMT\@. La operación de
intervalos
\[
 \mathcal Q(I,J)
 :=
 \left[
 \frac{\inf I}{(\sup J)^2},
 \frac{\sup I}{(\inf J)^2}
 \right]
\]
es el menor intervalo que contiene \(x/y^2\) para todo
\((x,y)\in I\times J\). Preserva el encaje: si los cilindros de entrada se
refinan, también se refina \(\mathcal Q(I_n^\pi,I_n^\varphi)\). Como sus
diámetros tienden a cero y el límite de \(I_n^\varphi\) es positivo,
\[
 \bigcap_n\mathcal Q(I_n^\pi,I_n^\varphi)
 =
 \left\{\frac{\pi_{\rm HMT}}{\varphi_{\rm HMT}^2}\right\}.
\]
Con los doce bloques en base mil ya publicados, el intervalo cociente fuerza
treinta y cinco cifras decimales:
\[
 \frac{\pi_{\rm HMT}}{\varphi_{\rm HMT}^2}
 =
 1.19998161486432666111577775331680692\ldots.
\]
Esta construcción conserva la procedencia de cada cifra: el numerador
procede del cilindro de clausura y el denominador de la autoescala.

\paragraph{Separación algebraica entre lectura areal, lectura volumétrica y sus codominios respectivos}
La razón completa no puede salir de la sola álgebra racional de
\(F_{\rm av}\). En efecto, sus funciones racionales espectrales pertenecen a
\(\mathbb Q(\sqrt5)\); allí se encuentra \(\varphi^{-2}\), pero no el número
trascendente \(\pi\). Por consiguiente, toda derivación que pretendiese
obtener \(\pi/\varphi^2\) únicamente de la matriz \(2\times2\) estaría
mal tipada. El dato trascendente debe proceder del lector coinductivo de
clausura, o de un cociclo de frontera equivalente demostrado
independientemente. Este control negativo fija la división de trabajo entre
la correspondencia finita y el lector de número.

\subsubsection{Involución de intercambio del cociente}
\label{sec:espejo-pi-phi-electron-hmt}

La correspondencia admite dos orientaciones exactas de una misma ley. Si
\[
 \mathscr R(x,y)=\frac{x}{y^2},
 \qquad
 \mathsf J(x,y)=(y,x),
\]
entonces \(\mathsf J^2=I\) y
\[
 \mathscr R(\pi,\varphi)=\frac{\pi}{\varphi^2},
 \qquad
 (\mathscr R\circ\mathsf J)(\pi,\varphi)
 =\frac{\varphi}{\pi^2}.
\]
Se verifican además las identidades
\[
 \mathscr R(\pi,\varphi)\mathscr R(\varphi,\pi)
 =\frac1{\pi\varphi},
 \qquad
 \frac{\mathscr R(\pi,\varphi)}
 {\mathscr R(\varphi,\pi)}
 =\left(\frac{\pi}{\varphi}\right)^3.
\]
Escribiendo
\(\mathcal B_{\rm av}^{(\pi)}=\mathscr R(\pi,\varphi)\), se obtiene
también la reparametrización exacta
\[
 \frac{\varphi}{\pi^2}
 =
 \frac1{\varphi^3
 \bigl(\mathcal B_{\rm av}^{(\pi)}\bigr)^2},
 \qquad
 e^{\varphi/\pi^2}
 =
 \exp\!\left[
 \frac1{\varphi^3
 \bigl(\mathcal B_{\rm av}^{(\pi)}\bigr)^2}
 \right].
\]
Así, el cambio de orientación entre cierre y crecimiento transporta la
razón de frontera a su lectura intercambiada:
\[
 \frac{\pi}{\varphi^2}
 \xleftrightarrow{\ \mathsf J\ }
 \frac{\varphi}{\pi^2}.
\]
La primera orientación lee el cierre areal marcado sobre la memoria
volumétrica cuadrática; la segunda lee el crecimiento interior sobre el
cierre cuadrático. Esta correspondencia es un teorema del lector HMT. La
Mecánica Dimensional aplica después la segunda orientación dentro de la
construcción electrónica, sin convertir al electrón en antecedente de
\(\pi\), \(\varphi\) o de la correspondencia.

\subsubsection{Principio de Ambivalencia de los lectores areal y volumétrico}
\label{sec:ambivalencia-grav-inercial-pi-phi2}

El Principio de Ambivalencia afirma que la lectura areal y la lectura
volumétrica son dos funcionales coordinados, pero no intercambiables, de un
mismo estado con memoria. Su coincidencia en la diagonal es
\[
 \mathcal L_{\rm ar}^{\rm diag}(x)
 =\mathcal L_{\rm vol}^{\rm diag}(x).
\]
La igualdad diagonal no borra la ambivalencia fuera de ella.

Sobre la frontera areal--volumétrica, ambos lectores factorizan a través de
una misma amplitud positiva de memoria \(\mathcal M(x)\):
\[
 \mathcal L_{\rm vol}(x)=\mu_{\rm vol}\mathcal M(x),
 \qquad
 \mathcal L_{\rm ar}(x)=
 \pi_{\rm HMT}\mu_{\rm ar}\mathcal M(x).
\]
Entonces el teorema de Parry y la marca de retorno dan
\[
 \boxed{\displaystyle
 \frac{\mathcal L_{\rm ar}(x)}
 {\mathcal L_{\rm vol}(x)}
 =\frac{\pi_{\rm HMT}}{\varphi^2}.}
\]
Con la normalización volumétrica, la diferencia relativa de lectores es
\[
 \frac{\mathcal L_{\rm ar}-\mathcal L_{\rm vol}}
 {\mathcal L_{\rm vol}}
 =
 \frac{\pi_{\rm HMT}}{\varphi^2}-1
 =
 0.1999816148643266611\ldots.
\]
La lectura física de esta correspondencia se realiza más tarde, en el capítulo de
Cartan--Holst, después de construir holonomía, co-tétrada, conexión, torsión
y curvatura. Aquí el resultado es enteramente HMT: una razón exacta entre
dos lectores de la misma memoria de hoja.

Conviene distinguirlos del extractor enriquecido que se utiliza después en
Mecánica Dimensional. Allí una historia completa conserva el conteo de acción
\(n_A\), los doce eventos orientados \(e_0,\ldots,e_{11}\) y la cocadena
torsional a nivel de paso; de ellos se obtienen
\[
 s_C=\sum_{i=0}^{5}(e_i+e_{i+6}),
 \qquad
 W_\pm=6n_A\pm s_C.
\]
Los canales ya generados por HMT satisfacen entonces la identidad exacta
\[
 \exp\!\left(n_AA_{\rm rad}+\frac{s_C}{6}C^\ast_{\rm rad}\right)
 =\left(q_+^{-W_+}q_-^{-W_-}\right)^{1/12}.
\]
Por tanto, el paso ruta--carácter angular no está ausente: se realiza sobre
el registro dodecafásico enriquecido. Lo que este cálculo de fase prueba por sí
solo son \(\Omega_{\rm sw},S,G\); no los sustituye ni permite reconstruir el
registro completo después de haberlo proyectado.

La medida uniforme sobre las \(104\,976\) presentaciones desciende a átomos
de masas \(1/729\), \(1/243\) y \(1/54\), pero no a \(1/468\). En la
microfibra quíntuple de \(\pi\), las dos probabilidades son
\[
 \tau_{468}(\mathscr R_\pi)=\frac5{468},
 \qquad
 F_\#\mu_{\rm seed}(\mathscr R_\pi)=\frac7{729}.
\]
Tanto la elevación promediada como la elevación de fase realizan la primera: no
empujan la medida uniforme sobre semillas, sino que asignan la misma masa
total a cada clase observable. La elevación de fase conserva además la
quietud microscópica en los pasos neutros.

Para relacionar la elevación de fase con las historias enriquecidas de
profundidad seis hace falta un núcleo \(\mathcal K_6\) sobre
\(\mathcal H_6^{\rm cat}\times C_9\). Éste desciende a
\(\widehat{\mathcal K}_{\rm ph}\) a través de
\(\rho_6\times\mathrm{id}_{C_9}\) si y sólo si satisface, en cada fase, la
condición de agregabilidad fuerte
\[
 \sum_{\rho_6(g)=y}\mathcal K_6((h,r),(g,r+1))
 =L_{\tau(r)}(\rho_6(h),y),
\]
independientemente de \(h\) dentro de cada fibra de \(\rho_6\). En ese caso,
la composición con \(F\) proyecta a \(\mathcal K_{\rm ph}\). La construcción
de una familia \((\mathcal K_m)_{m\geq6}\) que satisfaga esta identidad y
conmute con todos los truncamientos es el paso restante hacia una
transferencia del límite inverso.

\section{Cociclo de escala y conservación cilíndrica}
\label{sec:cociclo-parry-escala}

Sea \(r=(1,\varphi)\) un vector de Perron--Frobenius de
\(F_{\rm av}\). La matriz de transición de Parry es
\[
P_{ij}=\frac{(F_{\rm av})_{ij}r_j}{\varphi r_i}
=
\begin{pmatrix}
0&1\\
\varphi^{-2}&\varphi^{-1}
\end{pmatrix}.
\]
Para una ruta admisible \(\gamma:i\to j\) de longitud \(n\),
\[
\mathbb P(\gamma\mid i)=\frac{r_j}{\varphi^nr_i},
\qquad
\boxed{\chi_P(\gamma)=\varphi^n\frac{r_i}{r_j}.}
\]
La telescopía prueba
\[
\chi_P(\gamma_2\circ\gamma_1)
=\chi_P(\gamma_2)\chi_P(\gamma_1).
\]
En una ruta cerrada, \(\chi_P=\varphi^n\); su cuadrado y su inverso
generan las combinaciones
\[
C_n=\frac{\chi_P^2+\chi_P^{-2}}2,\qquad
S_n=\frac{\chi_P^2-\chi_P^{-2}}2,
\]
que satisfacen
\[
x_{n+1}=3x_n-x_{n-1}.
\]

Como \(F_{\rm av}\) es simétrica, sea
\[
\ell=\frac{r}{1+\varphi^2},\qquad \ell^{\mathsf T}r=1.
\]
La medida de un cilindro es
\[
\mu[x_0\cdots x_n]=\ell_{x_0}r_{x_n}\varphi^{-n}.
\]
Por tanto,
\[
\boxed{
V_n(x)
=\frac{\mu[x_0]}{\mu[x_0\cdots x_n]}
=\varphi^n\frac{r_{x_0}}{r_{x_n}}
=\chi_P(x_0\cdots x_n).}
\]
Para una escala de dimensión tipada \(D>0\),
\[
a_n^{(D)}(x)=V_n(x)^{1/D}
\]
cumple la ley de conservación cilíndrica
\[
\boxed{
\bigl(a_n^{(D)}(x)\bigr)^D
\mu[x_0\cdots x_n]=\mu[x_0].}
\]

En rango espacial tres y sobre retornos cerrados, escribimos
\(a_n:=a_n^{(3)}\) y obtenemos
\[
\boxed{
\frac{a_9}{a_0}=\varphi^3,\qquad
\frac{a_{27}}{a_0}=\varphi^9,\qquad
\frac{a_{108}}{a_0}=\varphi^{36}.}
\]
Si \(A_k:=a_{9k}\) para \(k\geq1\), entonces, para \(k\geq2\),
\[
A_{k+1}-2A_k+A_{k-1}
=A_{k-1}(\varphi^3-1)^2>0.
\]
La desigualdad prueba convexidad en profundidad de retorno; una aceleración
física requiere además una carta de duración.

El cálculo primario de
\(\operatorname{Sym}^2(F_{\rm av})\) en
\cref{sec:retorno-marcado-pi-phi2} demuestra que su espectro es
\(\{\varphi^2,-1,\varphi^{-2}\}\). Conservamos aquí únicamente el corolario
para el cociclo de escala: sobre una ruta cerrada con \(V_n=\varphi^n\), el
modo estable obedece
\[
\boxed{M_n=\varphi^{-2n}=V_n^{-2}.}
\]
La identificación de estos tres autoespacios con las tres hojas del TRIT
exige un entrelazador; la coincidencia de cardinalidad no lo sustituye. El
corolario quedaría refutado si la ruta cerrada no satisficiera
\(V_n=\varphi^n\) o si el modo estable no tuviera autovalor
\(\varphi^{-2}\).

\section{Teorema de prolongación dinámica: hipótesis globales y dominio certificado}

Se han construido y separado explícitamente los objetos siguientes del
problema finito:
\[
 F,
 \qquad
 F_{\rm obs},
 \qquad
 \mathcal K_{\rm ph},
 \qquad
 \widehat{\mathcal K}_{\rm ph},
 \qquad
 K_{\rm cat},
 \qquad
 \tau_{468},
 \qquad
 F_{\rm av},
 \qquad
 \operatorname{Sym}^2(F_{\rm av}),
 \qquad
 D_\pi,
 \qquad
 \mathcal Q.
\]
Para \(F_{\rm obs}\) se ha probado el retorno determinista y se ha exhibido un
testigo explícito de no descenso producto. Para \(\mathcal K_{\rm ph}\) se
han probado la doble estocasticidad, la irreducibilidad, el período nueve, el
retorno nonádico uniforme y la unicidad de
\(\pi_{\rm ph}\). Para \(\widehat{\mathcal K}_{\rm ph}\) se han probado la
agregabilidad fuerte, la conservación microscópica de la quietud y la medida
estacionaria compensada. \(K_{\rm cat}\) ya no es un promedio sin
procedencia: es la sombra de un paso al promediar la fase uniforme.
El cociente de modos de TPK induce \(F_{\rm av}\); su elevación cuadrática
separa los modos \(\varphi^2,-1,\varphi^{-2}\). El operador \(D_\pi\)
marca una sola vez el retorno de frontera con la clausura generada, y el
lector intervalar \(\mathcal Q\) conserva el encaje de los cilindros de
\(\pi\) y \(\varphi\). Estos cuatro objetos construyen
\(\pi_{\rm HMT}/\varphi_{\rm HMT}^2\) sin confundir la frecuencia temporal,
la medida uniforme y la medida de Parry.

La fuerza de esta inducción es relativa y concreta. La palabra TPK fuerza los
pesos \(1/3\). La interpretación APP de todas las continuaciones compatibles,
junto con el conteo uniforme, fuerza \(E_+\) y \(E_\times\). La trayectoria
determinista \(F_{\rm obs}\) no fuerza esos operadores y, de hecho, no los
produce. Éste no es un defecto del resultado, sino la diferencia matemática
entre evolución de un camino y transferencia de un espacio de caminos.

\ifdefined\HMTInternalTraceability
\begin{criterion}[Elevación enriquecida de la reversión]
\label{pf84:SYN31-RHO-LIFT}
Se busca un subdominio \(D_\rho\subset X_{\rm TPK}^{\rm enr}\) y un mapa
\[
 \Theta_\rho:D_\rho\longrightarrow X_{\rm TPK}^{\rm enr}
\]
que eleve la reversión entera \(27\mapsto72\) del
\cref{int:pf84:SYN30-RHO-DESC,int:pf84:SYN30-CROSS-WITNESS}. Debe conmutar con los
lectores residual, entero y genealógico, cuantificar el cociente y el
acarreo, y preservar orientación, ruta, retorno y, cuando el transporte sea
invertible, monodromía, región y supervivencia.
Sólo si \(D_\rho\) es total e invariante bajo \(\Theta_\rho\) puede exigirse
\(\Theta_\rho^2=I\) y denominarse involución.

No existe todavía esta construcción ni una prueba de involutividad. El
programa queda falsado si no existe el mapa, si falla uno de los cuadrados
de lectura, si se pierde memoria sin cuantificarla o si
\(\Theta_\rho^2\ne I\) en un dominio donde se reclame involución.
\end{criterion}

\begin{criterion}[Elevación enriquecida del cuadrado]
\label{pf84:SYN31-SQ-LIFT}
Se busca un subdominio \(D_s\subset X_{\rm TPK}^{\rm enr}\) y un mapa
\[
 \Theta_s:D_s\longrightarrow X_{\rm TPK}^{\rm enr}
\]
que eleve \(s(n)=n^2\), en particular \(27\mapsto729\), y conmute con los
lectores residual, entero y genealógico de
\cref{int:pf84:SYN30-SQ-DESC,int:pf84:SYN30-CROSS-WITNESS}. Debe cuantificar
cociente y acarreo y preservar orientación, ruta, retorno y, cuando el
transporte sea invertible, monodromía, región y supervivencia. Este mapa no
es involutivo y no se le impone
\(\Theta_s^2=I\).

La prueba permanece ausente. La inexistencia de \(\Theta_s\), el fallo de
un cuadrado de lectura o la pérdida no controlada de cualquiera de las
coordenadas enriquecidas refutarían el programa. La ablación
\(81\mapsto72\) al cambiar \(q^\times(9,9)\) es un control independiente:
no es una rama de \(\Theta_\rho\) ni de \(\Theta_s\), y no prueba ninguno
de los dos mapas.
\end{criterion}
\fi

La conclusión finita es completa en su dominio: el retorno uniforme, la
medida estacionaria, la agregabilidad a profundidad seis y la separación
entre evolución de una ruta y transferencia del espacio de rutas quedan
demostrados por los objetos construidos. En particular,
\(\widehat{\mathcal K}_{\rm ph}\) no se identifica con la cronología de
\(F_{\rm obs}\), como muestra \cref{prop:retorno-fobs-identidad}.
